% Verbatim copy of /Users/alingalatan/workspaces/grothe/handoff/manuscript/legacy/ch25-parts/89g-conditional-gaussian-contact.tex (file dated 2026-10-25). Only this line was added.
% Rank-free original-response certificate, 22 October 2026.
\subsection{Radial unimodality and a rank-free response certificate}
\label{sec:conditional-response-contact-20261022}

\paragraph{Scope and status.}
The result below is \st{P,C}, not a new global endpoint.
It replaces the rank cutoff in one family of support tests by a
dimension-free projection criterion, and covers the full visible
first-chaos range. Its analytic input is radial unimodality of the
visible density, rather than a Gaussian scale-mixture representation.
All remainders may share coordinates with the selected test.

\begin{lemma}[\st{P}: visibly averaged concavity]
\label{lem:radial-visible-concavity}
Let $L\in\mathbb R^2$ have a bounded radial nonincreasing probability
density and finite first moment. Then
$\psi(u)=\E|L+\sqrt u|$ is increasing and concave on $[0,\infty)$.
In particular this holds for
$L=Z(a+b|Z|^2)$, with $a>0$, $b\ge0$ and $Z$ standard proper Gaussian.
\end{lemma}
\begin{proof}
Put $f(r)=\E|L+r|$ and $W(r)=\E|L+r|^{-1}$.
The convolution of two nonnegative radial decreasing functions is
radial decreasing: by layer cake it suffices that the intersection
area of two disks decreases with the distance between their centers.
Truncation applies this fact to $1/|x|$, so $W$ decreases.
The density bound makes its local singularity integrable.
In two real dimensions $\Delta|x|=1/|x|$ distributionally, with no
point mass at zero. The radial equation and $rf'(r)\to0$ at zero give
\[
 rf'(r)=\int_0^r sW(s)\,ds,\qquad
 \psi'(u)=\frac1{4u}\int_0^u W(\sqrt s)\,ds.
\]
The last expression is positive and decreasing. Gaussian smoothing
and passage to the limit prove the same statement for nonsmooth
densities. For the specified $L$, set $P(t)=t(a+bt)^2$. Its planar
density at radius $\sqrt{P(t)}$ is
\[
 \frac{e^{-t}}{\pi(a+bt)(a+3bt)},
\]
which is bounded and decreasing.
\end{proof}
This is a visible-averaged comparison. It does not assert the false
Gaussian domination at a fixed visible shift. The radial convolution
fact is standard; a short layer-cake proof is available at
\url{https://canizo.org/convolution/}.

\begin{lemma}[\st{P}: conditional covariance comparison]
\label{lem:conditional-covariance-contact}
Suppose $(J,\mathcal G)$ is independent of $L$ from
Lemma~\ref{lem:radial-visible-concavity}. Conditionally on
$\mathcal G$, let $J$ be a centered real two-dimensional Gaussian
with arbitrary covariance $\Sigma$. Singular and improper complex
Gaussians are allowed. If $J\in L^2$, then, for $d\ge0$,
\[
 \E|L+dJ|\le\E|L+d\norm J_2G|,
\]
where $G$ is independent standard proper Gaussian. The expectations
with $+J$ and $-J$ are equal.
\end{lemma}
\begin{proof}
Diagonalize $\Sigma$. For a standard real two-dimensional Gaussian,
its radius and direction are independent. Given the radius $R$, the
angular average of the squared output radius is
$R^2\operatorname{tr}\Sigma/2$. Concavity of $\psi$ gives
\[
 \E[|L+dJ|\mid\mathcal G]
 \le K(d^2\operatorname{tr}\Sigma),\qquad
 K(v)=\E|L+\sqrt vG|=\E_{E\sim\mathrm{Exp}(1)}\psi(vE).
\]
The function $K$ is concave. Jensen over $\mathcal G$ finishes the
proof since $\E\operatorname{tr}\Sigma=\norm J_2^2$.
Conditional centering and Gaussian symmetry give the sign equality.
\end{proof}

\begin{lemma}[\st{P}: exact projection-price transfer]
\label{lem:projection-price-dilution}
Use the actual responses $h,k$, $Ph=tZ$, $F=\epsilon Qk$ and
$V=\norm F_2^2$. Write $p=1/\epsilon$. Let $J$ be a unit remote
cubic satisfying Lemma~\ref{lem:conditional-covariance-contact}, and
suppose $|\langle J,F\rangle|^2\ge\eta V$, $0<\eta\le1$.
For every $\ell>2b\ge0$, put
\[
 L=Z\{\ell+b(|Z|^2-2)\},\qquad
 h_{\ell,b}(v)=2\E|L+\sqrt vG|.
\]
Then the actual response objective satisfies
\begin{equation}\label{eq:price-dilution-contact}
 \boxed{\mathcal N\le t^2-2\ell t+2pb^2+
   \sup_{v\ge0}\{h_{\ell,b}(v)-p\eta v\}.}
\end{equation}
For $V=0$ the same bound holds without choosing $J$.
\end{lemma}
\begin{proof}
Rotate $J$ so $r=\Re\langle J,F\rangle\ge0$ and put
$m=\Re\langle Z(|Z|^2-2),h\rangle$. The two disks $|h\pm k|\le1$
give the paired support inequality
\[
 \ell t+bm+pdr\le\tfrac12\E(|L+dJ|+|L-dJ|).
\]
Since $\norm{Qh}^2\ge m^2/2$, completing the square in $m$ and
dropping only $-\norm{Pk}^2$ gives
\[
 \mathcal N\le t^2-2\ell t+2pb^2+pV-2pdr+
          \E(|L+dJ|+|L-dJ|).
\]
Choose $d=r/\eta$, so $pV-2pdr\le-p\eta d^2$, and use the
conditional covariance comparison. If $V=0$, use $d=0$ directly.
\end{proof}
The factor $\eta$ is exact for proportional scale completion:
$(2\lambda-1/\eta)/\lambda^2\le\eta$, with equality at
$\lambda=1/\eta$. Thus the unknown remainder costs its energy,
not an $O(\sqrt{1-\eta})$ norm error. This spends the $Qh$ square
once; it must not be added to another bound that has already used it.

\begin{theorem}[\st{P,C}: rank-free, all-row response region]
\label{thm:conditional-response-contact}
Let $\epsilon=1891/2500$. For any remote first-chaos subspace
$\mathcal A\subset Z^\perp$, let $\mathcal B$ be its remote
orthogonal complement and set
\[
 \mathcal M(\mathcal A)=
 \mathcal H_{1,0}(\mathcal A)\mathcal H_{1,1}(\mathcal B)
 \oplus\mathcal H_{0,1}(\mathcal A)\mathcal H_{2,0}(\mathcal B).
\]
If
\begin{equation}\label{eq:one-leg-contact-share}
 \norm{\Pi_{\mathcal M(\mathcal A)}F}_2^2
          \ge\frac{499}{500}\norm F_2^2,
\end{equation}
then, for every $0\le t\le\sqrt\pi/2$,
\[
 \boxed{\mathcal N<\frac{2853}{4000}=0.71325<\frac{500}{701}.}
\]
Both dimensions, all quadratic ranks, all phases and amplitudes are
arbitrary. The orthogonal remainder may contain every visible sector
and depend on every test coordinate. More generally it is enough to
have any unit remote conditionally centered Gaussian cubic $J$ with
$|\langle J,F\rangle|^2\ge499V/500$.
\end{theorem}
\begin{proof}
Every element of $\mathcal M(\mathcal A)$ is
$\sum_iX_iY_i(\mathcal B)+\bar X_iW_i(\mathcal B)$, hence is
conditionally centered Gaussian given $\mathcal B$, independently of
the visible $Z$. Choose $J$ along the indicated orthogonal projection
and apply Lemma~\ref{lem:projection-price-dilution} with $\eta=499/500$.
The following six rational profiles cover the middle interval:
\[
\begin{array}{c|c|c|c}
 t\text{ interval}&\ell&b&S_{\ell,b}\text{ upper ceiling}\\\hline
 {[.697,.810406]}&.743838&.114290&1.2272692769557\\
 {[.810406,.835533]}&.816962&.134894&1.3322141699368\\
 {[.835533,.842503]}&.838731&.140930&1.3641895048025\\
 {[.842503,.847010]}&.844808&.142605&1.3731735402269\\
 {[.847010,.852507]}&.857148&.145993&1.3914924321810\\
 {[.852507,.880000]}&.867631&.148855&1.4071348626997
\end{array}
\]
Here $S_{\ell,b}=\sup_{v\ge0}(h_{\ell,b}(v)-p\eta v)$;
every terminating decimal is exact. The row cap in
\eqref{eq:price-dilution-contact} is convex in $t$. Exact comparisons
of its twelve endpoints give a maximum below $0.713249286$.

For $0\le t\le.697$, the $b=0$ support family gives
$t^2+p^2r^2\le\pi/4$. Therefore
\[
 \mathcal N\le\frac{\pi}{4p\eta}
       +\left(1-\frac1{p\eta}\right)t^2<0.712873.
\]
For $.88\le t\le\sqrt\pi/2$, take $\ell=1,b=1/6$.
Then $h'(0)=(3\pi/2)e^4\operatorname{erfc}(2)<p\eta$,
so concavity puts the scalar maximum at zero, with
$S_{1,1/6}=h(0)=11\sqrt\pi/12$.
The row quadratic decreases on this interval and is below $0.712597$.
\end{proof}

\paragraph{Directed scalar replay \st{C}.}
With $a=\ell-2b$ and $P(t)=t(a+bt)^2$, the two integrals are
\[
\begin{split}
 h(v)&=\sqrt{\pi v}\int_0^\infty e^{-t}
       {}_1F_1(-1/2;1;-P(t)/v)\,dt,\\
 h'(v)&=\frac{\sqrt\pi}{2\sqrt v}\int_0^\infty e^{-t}
       {}_1F_1(1/2;1;-P(t)/v)\,dt.
\end{split}
\]
For each row the fresh replay encloses $h(v_l),h'(v_l),h'(v_r)$
and verifies $h'(v_l)>p\eta>h'(v_r)$. The tangent at $v_l$ then
bounds the supremum by
\[
 h(v_l)-p\eta v_l+(h'(v_l)-p\eta)(v_r-v_l).
\]
Complex ball integration on $[0,40]$ has an entire integrand.
The two respective positive tail bounds are
\[
 e^{-40}\{2(41a+1682b)+\sqrt{\pi v}\},\qquad
 e^{-40}\frac{\sqrt\pi}{2\sqrt v}.
\]
The first uses the triangle inequality and $\sqrt t\le t$ for
$t\ge40$; the second uses the beta integral to bound
${}_1F_1(1/2;1;-x)\le1$. Two isolated-directory replays at 110 and
160 bits pass. The separate standard-library audit checks the two
outer intervals, all rational endpoints, and the elementary algebra.
It assumes the six integral ceilings unless the full replay is used.
The package is
\pathref{handoff/research-third-decimal/conditional-response-contact-20261022/};
\code{verify\_all.py} requires \code{python-flint==0.9.0}.
No numerical assertion here is a floating-point acceptance test.

\paragraph{An exact invariant entrance condition \st{P,C}.}
Let $q=\E[F\mid Z^\perp]$ be the remote cubic. Let $N_{\mathcal A}$
count both holomorphic and antiholomorphic $\mathcal A$ legs, and set
$I_j=\langle q,(N_{\mathcal A})_{\underline j}q\rangle$, $j=1,2,3$.
These are the ordered $j$-fold derivative-energy sums, including both
Wirtinger derivative types and repeated coordinate indices. Since
$N_{\mathcal A}$ has spectrum $0,1,2,3$ on cubics,
\begin{equation}\label{eq:one-leg-projection}
 \boxed{\norm{\Pi_{\mathcal M(\mathcal A)}q}^2=I_1-I_2+I_3/2.}
\end{equation}
Indeed $n-n(n-1)+n(n-1)(n-2)/2$ equals $1_{n=1}$ on these four
integers. Thus $I_1-I_2\ge499V/500$ is a simpler sufficient test.
The condition uses the remote projection, not a visible-dependent
test disguised as remote.

For the coordinate-core obstruction in
Theorem~\ref{thm:joint-radial-core},
\[
 Q_N=\frac{X_1(\sum_{j=1}^N|X_j|^2-N-1)}{\sqrt{N+1}},
\]
take $\mathcal A=\operatorname{span}(X_1)$. The one-leg share is
$(N-1)/(N+1)$, and the self-interacting cubic remainder has share
$2/(N+1)$. Consequently every actual response with $F=cQ_N$ is
covered for $N\ge999$. Its arbitrary coupled neighbourhood given by
\eqref{eq:one-leg-contact-share} is covered too. The fixed-core
approximation obstruction remains true, but no longer obstructs this
response branch.

\paragraph{The precise remaining gap.}
The new region is not exhaustive. A unit diagonal Wick cube has largest
one-leg share $4/9$: for $r=\norm{\Pi_{\mathcal A}X}^2$ its share is
$3r(1-r)^2$, and
$4/9-3r(1-r)^2=(r-1/3)^2(4-3r)$. Such cubes use the older Wick
contact cone. No theorem forces a general surviving cubic into the
union of those cones and \eqref{eq:one-leg-contact-share}.
In particular, failure of the latter does not imply small influence,
independent blocks, or a small coordinate core. The new branch has a
complete all-row cover at the target witness; its coupled complement
still needs one. The analytic proof and ball enclosures are not yet
Lean theorems.
